\chapter{الهندسة الاقليدية}

\section[مقدمة]{مقدمة في الهندسة الإقليدية \en{(Euclidean Geometry)} \cite{ref1}}
بدأ كل شيء في القرن الثالث قبل الميلاد، عندما قام العالم اليوناني إقليدس في كتابه "الأصول" بصياغة نظام رياضي متكامل يعتمد على مجموعة من البديهيات والمسلمات التي تبدو بديهية للعقل البشري. اعتمدت هذه الهندسة على فكرة "المستوي" المسطح، حيث الخط المستقيم هو أقصر مسافة بين نقطتين، ومجموع زوايا المثلث يساوي دائماً 180°.
كان النظام الإقليدي يُعتبر الحقيقة الوحيدة الممكنة لوصف الكون، حتى جاءت المسلمة الخامسة (مسلمة التوازي) لتثير تساؤلات العلماء. تنص هذه المسلمة على أنه من نقطة خارج خط مستقيم، لا يمكن رسم إلا خط واحد فقط يوازي الخط الأصلي.
لفترة تزيد عن ألفي عام، حاول الرياضيون إثبات أن المسلمة الخامسة ليست "بديهية" بل هي "مبرهنة" يمكن استنتاجها من المسلمات الأربع الأولى، لكن جميع المحاولات باءت بالفشل. هذا الفشل أدى في النهاية إلى اكتشاف مذهل في القرن التاسع عشر على يد علماء مثل لوباتشيفسكي، بولياني، وريمان.
بدأ العلماء بتجربة "افتراض نفي" المسلمة الخامسة لرؤية ما إذا كان ذلك سيؤدي إلى تناقض منطقي. بدلاً من التناقض، اكتشفوا عوالم هندسية جديدة ومتسقة تماماً:
1. الهندسة الزائدية (Hyperbolic Geometry): حيث افترضوا أنه يمكن رسم عدد لانهائي من الخطوط الموازية من نقطة واحدة. هنا، يظهر المكان بشكل "سرج الفرس"، ويكون مجموع زوايا المثلث أقل من 180°.
2. الهندسة الناقصية (Elliptic Geometry): حيث افترضوا عدم وجود خطوط متوازية على الإطلاق. في هذا النظام (مثل سطح الكرة)، يلتقي أي خطين مستقيمين في النهاية، ويكون مجموع زوايا المثلث أكبر من 180°.
لم يكن التحول من الإقليدية إلى اللاإقليدية مجرد ترف فكري، بل كان تغيراً جذرياً في فهمنا للكون. فقد أدركنا أن هندسة إقليدس هي حالة خاصة تنطبق فقط على الأسطح المسطحة

\section[لمسافة بين نقطتين في المستوى الإقليدي]{المسافة بين نقطتين في المستوى الإقليدي \cite{ref1}}

في الهندسة الإقليدية، تُعرف المسافة بأنها طول القطعة المستقيمة الواصلة بين نقطتين. ويتم اشتقاق قانون المسافة مباشرة من نظرية فيثاغورس.\\
\textbf{الاشتقاق الرياضي:}\\
إذا كان لدينا نقطتان في المستوى الإحداثي $A(x_1, y_1)$ و $B(x_2, y_2)$، فإننا نتخيل مثلثاً قائم الزاوية يكون فيه الفرق في السينات $(\Delta x)$ هو القاعدة، والفرق في الصادات $(\Delta y)$ هو الارتفاع.\\
\textbf{القانون:}\\
تُحسب المسافة $d$ باستخدام الصيغة التالية:
\[
d = \sqrt{(x_2 - x_1)^2 + (y_2 - y_1)^2}
\]
\begin{example}
	جد المسافة بين النقطتين
	$A(1, 2)$ و $B(4, 6)$
\end{example}
\begin{solution}
	المسافة الإقليدية بين النقطتين $A(1,2)$ و $B(4,6)$ تُعطى بالعلاقة:
	\[
	d(A,B)=\sqrt{(x_2-x_1)^2+(y_2-y_1)^2}
	\]
	بالتعويض:
	\[
	d(A,B)=\sqrt{(4-1)^2+(6-2)^2}
	\]
	\[
	d(A,B)=\sqrt{3^2+4^2}
	\]
	\[
	d(A,B)=\sqrt{9+16}
	\]
	\[
	d(A,B)=\sqrt{25}
	\]
	\[
	d(A,B)=5
	\]
\end{solution}
\section[معادلة المستقيم في الهندسة الإقليدية]{معادلة المستقيم في الهندسة الإقليدية \cite{ref2}}
المستقيم في المنظور الإقليدي هو "طول بلا عرض" يمتد إلى ما لا نهاية في جهتين متضادتين، ويتميز بأن له ميلاً ثابتاً (Slope) لا يتغير بين أي نقطتين عليه.
\\\begin{tasks*}
	\item ميل المستقيم $m$:\\
	هو نسبة التغير الرأسي إلى التغير الأفقي:
	\[
	m = \frac{y_2 - y_1}{x_2 - x_1}
	\]
	\item الصيغ العامة لمعادلة المستقيم\\
		يمكن تمثيل المستقيم في المستوى بعدة صيغ مكافئة، أهمها:
	
	\begin{enumerate}
		\item \textbf{الصيغة العامة:}
		\[
		Ax + By + C = 0
		\]
		
		\item \textbf{صيغة الميل والمقطع:}
		\[
		y = mx + b
		\]
		حيث $m$ هو الميل و $b$ هو نقطة التقاطع مع محور $y$.
	\end{enumerate}
\end{tasks*}

\begin{example}
	أوجد معادلة المستقيم المار بالنقطتين:
	\[
	A(1,2),\quad B(4,6)
	\]
\end{example}
\begin{solution}
	نحسب الميل:
	\[
	m=\frac{6-2}{4-1}=\frac{4}{3}
	\]
	نستخدم صيغة:
	\[
	y=mx+b
	\]
	نعوض بالنقطة $A(1,2)$:
	\[
	2=\frac{4}{3}+b
	\]
	\[
	b=\frac{2}{3}
	\]
	إذن:
	\[
	y=\frac{4}{3}x+\frac{2}{3}
	\]
	نحوّل للصيغة العامة:
	\[
	3y=4x+2
	\]
	\[
	4x-3y+2=0
	\]
\end{solution}
\renewcommand{\thetable}{}
\begin{table}[ht!]
	\centering
	\renewcommand{\arraystretch}{1.5}
    \begin{tabular}{|p{4.1cm}|p{4.1cm}|p{4.1cm}|p{4.1cm}|}
		\hline
		{}& \textbf{اقليدية} & \textbf{هذلولية} & \textbf{اهليجية} \\\hline
		يتقاطع المستقيمان في & نقطة واحدة في الأكثر & نقطة واحدة في الأكثر & نقطة واحدة (احادية) ، نقطتان (المزدوجة)\\\hline
		ليكن $l$ مستقيماً و $p$ نقطة لا تقع على $l$ فأنه & يوجد موازي واحد وواحد فقط للمستقيم $l$ من النقطة $p$ & يوجد موازيان للمستقيم $l$ من النقطة $p$ & لا يوجد موازي الى $l$ من $p$ \\\hline
		يفصل المستقيم الى نصفين & بواسطة نقطة & بواسطة نقطة & كلا\\\hline
		اذا قطع مستقيم احد مستقيمين متوازيين فأنه & يقطع الآخر & قد يقطع وقد لا يقطع & ------------\\
		\hline
		العمودان على نفس المستقيم يكونان & متوازيين & غير متوازيين & متقاطعان \\
		\hline
		مجموع زوايا المثلث  & تساوي زاويتين قائمتين & اقل من زاويتين قائمتين & اكثر من زاويتين قائمتين\\
		\hline
	\end{tabular}
	\caption{مقارنة بين الهندسة الإقليدية واللاإقليدية}
\end{table}